\documentclass[aspectratio=169]{beamer}

\usetheme{Madrid}
\usepackage{booktabs}
\usepackage{xcolor}
\setbeamertemplate{navigation symbols}{}

\title[Review fixes, any-node models, pan-India eval]{Acting on the review: provenance, any-node models, and honest pan-India numbers}
\author[Salil Gujar]{Salil Gujar \\ \footnotesize M.Tech CSE \quad Entry No.\ 2025MCS2106}
\date{August 2026}

\begin{document}

\frame{\titlepage}

% ------------------------------------------------------------------
\begin{frame}{Where we are, and what this tackled}
\begin{columns}[T]
\begin{column}{0.5\textwidth}
\emph{Done so far}
\begin{itemize}\small
  \item A 4-class base map of India at 10\,m with a living hierarchy: split, add, rule-split, merge,
  retrain, all as crisp Earth-Engine tiles.
  \item A git-backed zoo of cards, save and resume, the lab's IndiaSAT models plugged in, per-fortnight
  water, STACD provenance.
\end{itemize}
\end{column}
\begin{column}{0.5\textwidth}
\emph{This round: acting on the review}
\begin{itemize}\small
  \item Cross-check the STACD record with the other teams and align it.
  \item Train the lab's model the way they do.
  \item Let a user attach a model to any node, not only greenery.
  \item Measure mining and water honestly, pan-India, on ground truth.
  \item Add a filter for spurious water; improve acacia and mining.
\end{itemize}
\end{column}
\end{columns}
\end{frame}

% ------------------------------------------------------------------
\begin{frame}{STACD: cross-checking with the other teams}
Susmit checked our record against theirs, flagged that some parameters differ, and was unsure which are
optional or mandatory. We compared field by field and aligned ours.
\begin{itemize}\small
  \item The mandatory STAC fields already matched. The differences were optional metadata we now adopt
  for uniformity: the version, a collection, real catalog links, a datetime range, keywords, and the run
  parameters.
  \item Their pipeline-specific fields, a per-crown table description and drone-model names, do not apply
  to our raster output, which carries a class legend instead.
  \item A couple of format details remain open; following up with Saharsh this week, then re-sending.
\end{itemize}
\end{frame}

% ------------------------------------------------------------------
\begin{frame}{Training the lab's model the way they do}
A correctness check on the farm, plantation, scrubland model we ported.
\begin{itemize}\small
  \item We were training it only on ground truth within forty kilometres of the user's box, so a small
  box saw a weaker model than the one the lab ships.
  \item The lab trains it across the whole agro-ecological region. We now do the same: drop the
  box-local restriction, sample the whole region with a balanced cap per class.
  \item The tree-against-crop model already trained on its full pan-India asset, so it was left alone.
  \item Why it is re-trained each time and not stored: it is an Earth-Engine model that trains and runs
  on Google's servers, so there is no file to save; the same fixed data and seed make it reproducible.
\end{itemize}
\end{frame}

% ------------------------------------------------------------------
\begin{frame}{Any model on any node, with a suggestion}
A core promise is that a user refines any class, not a fixed one.
\begin{itemize}\small
  \item The lab's two models were hard-wired to refine greenery. They can now attach to any node the
  user selects.
  \item The card still suggests where it normally goes, greenery, but does not decide for the user.
  \item Worked example: split greenery by a rule, then refine one of those children with the lab's
  model; only that child changes, the rest of the map stays.
  \item A guard refuses a clash cleanly, and the machinery underneath already generalised, so this was a
  small, safe change.
\end{itemize}
\end{frame}

% ------------------------------------------------------------------
\begin{frame}{Is the mining pixel-to-polygon good enough?}
The question: does vectorising the mining pixels avoid a heavier, learned segmentation model? Measured
pan-India, outside the framework, on real mine polygons.
\begin{center}
\begin{tabular}{lrrr}
\toprule
Mining, on ground-truth polygons & Precision & Recall & F1 \\
\midrule
Vectorised polygons vs truth (object level) & 0.04 & 0.16 & 0.07 \\
Pixel classifier, linear & 0.45 & 0.70 & 0.55 \\
Pixel classifier, forest with tuned cut & 0.61 & 0.58 & 0.59 \\
\bottomrule
\end{tabular}
\end{center}
\begin{itemize}\small
  \item The vectorised polygons over-fragment and rarely line up with a real mine, so they are not a
  delineator; a learned segmentation model is the route for true objects.
  \item The pixel classifier is a usable screen: a forest with a lowered decision cut lifts precision
  from 0.45 to 0.61, the weak spot, at some cost to recall.
\end{itemize}
\end{frame}

% ------------------------------------------------------------------
\begin{frame}{Water: small against large}
The water model currently runs on Sentinel-1 radar, with Sentinel-2 optical water indices alongside.
Measured on the deployed model, split by the size of the water body.
\begin{center}
\begin{tabular}{lrrr}
\toprule
Deployed model, by water-body size & Precision & Recall & F1 \\
\midrule
Large bodies & 1.00 & 0.97 & 0.99 \\
Small bodies & 0.85 & 0.67 & 0.75 \\
\bottomrule
\end{tabular}
\end{center}
\begin{itemize}\small
  \item Large water is essentially solved.
  \item Small water is the real gap: a third of it is missed, a recall problem, not a spurious one, since
  dry land is wrongly called water only about two percent of the time.
\end{itemize}
\end{frame}

% ------------------------------------------------------------------
\begin{frame}{A correction for spurious water, checked on ground truth}
The review asked for a threshold that holds water only where it persists, a correction on the water
output. We fold it into how annual water is decided and validate it on the lab's Earth-Engine water
ground truth.
\begin{center}
\begin{tabular}{lrrr}
\toprule
Persistence threshold & Water precision & Spurious rate & Small-water recall \\
\midrule
Any fortnight & 0.80 & 0.15 & 0.30 \\
At least two fortnights & 0.96 & 0.02 & 0.11 \\
\bottomrule
\end{tabular}
\end{center}
\begin{itemize}\small
  \item Holding a pixel as water only if it was wet in at least two fortnights cuts spurious water from
  fifteen percent to two, and lifts precision to 0.96.
  \item The cost lands on small and seasonal water, so a single global threshold is blunt: it argues for
  the two-classifier design, a lenient first pass that lets seasonal water through, then a stricter one
  inside water bodies.
\end{itemize}
\end{frame}

% ------------------------------------------------------------------
\begin{frame}{Acacia: how many crowns, and a fair filter}
\begin{itemize}\small
  \item We have 336 acacia crowns and 576 non-acacia. Every crown is a single tree, a median of
  twenty-seven square metres.
  \item That is smaller than one ten-metre pixel, so each crown is a mixed pixel, tree blended with the
  ground around it. This is why the split is hard on this embedding.
  \item The review's ten-by-ten filter, taken literally, drops ninety-eight percent of the data. Instead
  we drop only degenerate slivers under fifteen square metres and keep the rest, 296 and 498 crowns.
  \item So the filter is really a diagnosis: our acacia ground truth is thin and sub-pixel.
\end{itemize}
\end{frame}

% ------------------------------------------------------------------
\begin{frame}{Acacia: improving it, and the ceiling}
\begin{center}
\begin{tabular}{lrrr}
\toprule
Configuration & Precision & Recall & F1 \\
\midrule
Linear, single year (baseline) & 0.58 & 0.82 & 0.68 \\
Linear, several years pooled & 0.59 & 0.84 & 0.69 \\
Forest, several years pooled & 0.68 & 0.74 & 0.71 \\
\bottomrule
\end{tabular}
\end{center}
\begin{itemize}\small
  \item The proven levers, more years and a non-linear model, lift F1 from 0.68 to 0.71 and accuracy from
  0.72 to 0.78, mostly by raising precision. Real, but bounded: the input is still a mixed pixel.
  \item The true ceiling-raiser is higher-resolution features, Tessera or drone imagery encoded with
  DINO, which is external. Nothing here is solved; it is honest and improved.
\end{itemize}
\end{frame}

% ------------------------------------------------------------------
\begin{frame}
\centering
\vspace{1.5cm}
{\Large Thank you}
\vspace{0.5cm}


\end{frame}

\end{document}
